\documentclass[11pt]{article}
\usepackage[utf8]{inputenc}
\usepackage[T1]{fontenc}
\usepackage{graphicx}
\usepackage{booktabs}
\usepackage{amsmath}
\usepackage{amssymb}
\usepackage[margin=1in]{geometry}
\usepackage{caption}
\usepackage{array}
\begin{document}
\title{Mothwick compress and symptom clearance: a review of small parallel comparisons}
\author{Grove Room Review Group}
\maketitle
\begin{abstract}
Small reports compared a Mothwick compress with plain wrap for clearance of brindle ache.
{{v:n_included}} studies remained in the quantitative summary.
The pooled risk difference, from fixed-effect inverse-variance weighting, was {{v:rd_pool}} percentage points
(95\% CI {{v:rd_lo}} to {{v:rd_hi}}).
Heterogeneity was Q = {{v:q_stat}}, df = {{v:q_df}}, and $I^{2} = {{v:i2_pct}}\%$.
\end{abstract}
\section{Introduction}
Brindle ache is a short-lived burning discomfort. A cloth compress has been compared with plain wrap in several small parallel groups, with clearance counted on a shared day \cite{quill2021,brindle2015}.
Northline, Keelhorn, Orrin, Pavelic, Suttor, and Wynnell each reported that comparison \cite{northline2014,keelhorn2016,orrin2017,pavelic2018,suttor2019,wynnell2020}.
This review pools those reports. It does not add unpublished counts.
\section{Methods}
Reports were eligible when both groups were followed to the clearance day and the event counts were printed for each group.
Clearance was the absence of brindle ache on day {{v:day_clear}}.
The effect is the risk difference in percentage points, one hundred times the treated proportion minus the plain-wrap proportion.
Each study interval is a 95\% Wald interval with multiplier 1.96.
The summary is a fixed-effect inverse-variance pool.
The weight of a study is one divided by the square of its standard error, and the printed weight is that quantity as a percentage of the total weight.
The pooled risk difference is the weighted mean of the study risk differences.
Its 95\% interval uses the multiplier 1.96 on one divided by the square root of the sum of the weights.
Cochran's Q is the weighted sum of squared deviations from the pooled estimate, with degrees of freedom equal to the number of studies minus one.
The $I^{2}$ percentage is the larger of zero and one hundred times (Q minus those degrees of freedom) divided by Q.
Risk differences, weights, interval endpoints, Q, and $I^{2}$ are rounded half up to one decimal place \cite{brindle2015}.
\section{Results}
Records identified through database searching: {{v:n_identified}}.
Duplicates removed: {{v:n_duplicates}}.
Records screened: {{v:n_screened}}.
Records excluded: {{v:n_rec_excl}}.
Full-text articles assessed for eligibility: {{v:n_fulltext}}.
Full-text articles excluded: {{v:n_ft_excl}}.
Studies included in the meta-analysis: {{v:n_included}}.
The full-text reports were excluded because the population did not match (n = {{v:n_ex_pop}}), the comparator was not plain wrap (n = {{v:n_ex_comp}}), clearance on day {{v:day_clear}} was not reported (n = {{v:n_ex_out}}), or the same cohort had already been reported (n = {{v:n_ex_cohort}}).
\falsifytable{T1}
The effect table sums to {{v:e_tx}} of {{v:n_tx}} with the compress and {{v:e_ctrl}} of {{v:n_ctrl}} with plain wrap.
The inverse-variance pooled risk difference was {{v:rd_pool}} percentage points (95\% CI {{v:rd_lo}} to {{v:rd_hi}}).
Study risk differences ran from {{v:rd_s5}} to {{v:rd_s4}} percentage points.
The analysis weights sum to {{v:w_sum}}.
Heterogeneity was Q = {{v:q_stat}}, df = {{v:q_df}}, and $I^{2} = {{v:i2_pct}}\%$.
\falsifytable{T2}
\falsifyfigure{F1}
\section{Discussion}
{{v:ph_concl_claim}}
{{v:ph_hedge}}
Several study intervals included a null difference, so the pooled estimate should not be read as a precise ranking of the reports \cite{brindle2015}.
The clearance wording follows one definition, and a report that counted a different day would not have entered the table \cite{quill2021}.
Limitations.
{{v:ph_limit}}
\falsifybibliography
\end{document}
